% REVIEW DRAFT: not an input of the current published manuscript.
% This file is not proof certification; see geometry-audit-2026-11-21.json.
\section{The integral-point eigenlattice comparison}
\label{sec:sr-integral-eigenlattice-comparison}

We specify the integral-model hypotheses of the comparison.
Let \(E\) be the CM field of the PEL datum, let \(S\)
be a finite set of rational primes containing every prime
dividing \(\operatorname{disc}E\), and let
\(\mathfrak M\) be the normal Hodge compactification used
for the height comparison. Its semiabelian group is smooth;
its generic fibre is the given principally polarized
abelian scheme with \(\mathcal O_E\)-action.
At \(p\in S\) its chosen torsion level is \(p^{R_p}\),
where \(R_p\ge4\), and full generic torsion and the
polarization's Weil pairing are part of its data.
Assume the precise Hodge--Tate freeness condition:
the finite-level cotangent lattice \(\omega'_p\)
and, for every CM label \(s\), its image in the dual
of the labelled generic Lie space are locally free
on the \(p\)-fibre, of their prescribed ranks.
These are both clauses of \(\mathrm{HTFree}(\mathfrak M,p)\).

Let \(z:\operatorname{Spec}\mathcal O_L\to\mathfrak M\)
be the actual integral point.
Transport the Lie space of its given abelian-part model
using the cartesian Lie identification in the point data.
Write \(V\) for this \(L\)-space and
\(\Lambda\subset V\) for the integral tangent module
of the smooth group along \(z\).
The given \(\mathrm{LinPt}\) identifies this with its
integral tangent lattice, its action with the model's
\(\mathcal O_E\)-action, and its labelled field embeddings
with the model's labelled functions.
For a label \(s\), put
\[
 \Lambda_s=\Lambda\cap V_s,\qquad
 F_s=\mathfrak M.\mathrm{FibLat}_s(z).
 \tag{IC1}
\]
Here \(F_s\) is the intersection of the local saturated
evaluations of the patched Hodge lattice. The zero-prime term of this intersection puts \(F_s\)
in \(V_s\).
This definition, rather than a naive tensor product
of an intersection, is the one used below.

\paragraph{The different clears the integral projector.}
Let \(\sigma_s:E\hookrightarrow L\) be the labelled
embedding. Its restriction sends
\(\mathcal O_E\) into \(\mathcal O_L\), by integrality.
If \(a_i,b_i\) are trace-dual bases of \(E/\mathbb Q\),
with \(a_i\in\mathcal O_E\), the separability projector is
\[
                e_s=\sum_i\sigma_s(b_i)\rho(a_i).
 \tag{IC2}
\]
It acts as the identity on \(V_s\) and as zero on the
other eigenspaces. Since
\(b_i\in\mathfrak D_{E/\mathbb Q}^{-1}\), for every
\(d\in\mathfrak D_{E/\mathbb Q}\),
\[
                    \sigma_s(d)e_s\Lambda\subseteq\Lambda .
 \tag{IC3}
\]
Indeed \(db_i\in\mathcal O_E\), so the displayed
operator is an integral linear combination of
the \(\rho(a_i)\).
At primes where the different is a unit this is
an integral direct-sum decomposition.
At ramified primes (IC3) is used as stated; no
integral spectral idempotents are assumed there.

\paragraph{The local Hodge--Tate constant and its exact scope.}
For the finite-level Hodge lattice the constant is
\[
        \kappa_2=4,\qquad
        \kappa_p=\zeta_p-1\quad(p\ne2),
 \tag{IC4}
\]
where \(\zeta_p\) is primitive of order \(p\).
The actual chart construction, expanded below, gives
\[
       \omega'_p\subseteq\omega,\qquad
                   \kappa_p\omega\subseteq\omega'_p .
 \tag{IC5}
\]
The published input is the integral Hodge--Tate theorem
for a \(p\)-divisible group over a complete rank-one
valuation ring: its linearized map has cokernel killed
by an element of valuation at least \(1/(p-1)\);
its rational kernel is the orthogonal complement
specified by the dual Hodge--Tate map.
We use the conservative constant \(4\) at \(2\),
including the extra factor in the \(2\)-adic period
comparison.\footnote{L.~Fargues,
\href{https://www.numdam.org/item/10.24033/asens.2157.pdf}
{\emph{La filtration canonique des points de torsion des
groupes \(p\)-divisibles}}, Section 5.3.2, Theorem 2,
and the \(p=2\) qualification in Remark 2 immediately
before that section.}
This input applies to the actual divisible-group
test object constructed from the semiabelian torsion
tower, not directly to an arbitrary lattice called
Hodge--Tate.

We describe the match and the passage from test objects
to (IC5). On a normal affine chart \(A\) over
\(\mathbb Z_{(p)}\), with fraction field \(\Omega\),
the finite torsion group and Weil pairing give local
differential representatives \(h_x\), one for each
generic \(p^{R_p}\)-torsion point, with
\[
 \omega'_p=\sum_x A h_x+p^{R_p}\omega,\qquad
                 h'_x-h_x\in p^{R_p}\omega .
 \tag{IC6}
\]
The same formulas hold on overlapping charts.
The lower chart construction has the
\(\mathcal O_E\)-action, adjoint polarization pairing,
and the integral symplectic basis modulo \(p\).
The ranks are \(2g=\operatorname{rank}_{\mathbb Z}\mathcal L\)
and the prescribed label ranks.
They are the \(\mathrm{HTChart}\) and
\(\mathrm{IsHTFamily}\) data; the projectivity of
\(\omega'_p\) and its label images is exactly the
\(\mathrm{HTFree}\) hypothesis.

To test inclusion at a height-one prime above \(p\),
pass to a dominating complete algebraically closed
rank-one valued field.
The Tate parametrization of the generic abelian
scheme gives its compatible full torsion basis.
Its quasi-finite torsion models have a finite open
subgroup containing the closed fibre.
The corresponding subgroup \(H_n\subset\mathcal L/p^n\)
has the finite flat coordinate order.
Induce that order to all of \(\mathcal L/p^n\);
as a module it is the product of its translates over
the finitely many cosets of \(H_n\).
The subgroup identities
\(pa\in H_{n+1}\Longleftrightarrow a\in H_n\)
and the surjection \(H_{n+1}\to H_n\) give the
transition exact sequences.
These induced Hopf orders therefore define the
\(p\)-divisible test group.
Their cotangent module is the cotangent module of
the original smooth group: the identity finite open
part is unchanged.
The actual cubical polarization is used here.
Its translation isomorphisms and commutator construction
extend the Weil pairing as a regular function on
the product of the quasi-finite torsion schemes.
Explicitly, the ratio
\(\mathrm{cubN}(\theta,n,1,z,x)/
\mathrm{cubN}(\theta,n,1,x,z)\)
is a unit: the line-bundle factors cancel when
\(x^n=z^n=1\). This formula is functorial on every
test scheme and agrees over a proper geometric
fibre with the theta-group commutator;
this is \texttt{WeilPairingCubical} and
\texttt{TorsionTowerWeilPairing.weil\_regular}.
Restriction to the finite parts and the coset
induction preserve that regularity.
Together with the integral roots of unity and
the \(\mathcal O_E\)-action this supplies
\(\mathrm{PairingRegular}\), the polarized Tate
pairing and its adjoint action.
Regularity is not inferred from the generic
bilinear form on a possibly nonnormal finite group
scheme. The compatible PEL Tate parametrization
just described supplies the required integral
symplectic and \(\mathcal O_E\)-linear labels.
This is the finite-part construction used in
\(\mathrm{SemiabelianAdicTorsionOne}\)--%
\(\mathrm{Five}\); it does not assert that every
boundary torsion point extends as a section of the
semiabelian group.

The test group's finite-level differential is \(h_x\)
modulo \(p^{R_p}\omega\), by the Weil-character
definition. The integral Hodge--Tate theorem thus
puts \(\kappa_p\omega\) in the base-changed module
in (IC6).
The constants satisfy
\(\kappa_p\mid p^2\), so also
\(\kappa_pp^{R_p-2}\mid p^{R_p}\).
Projectivity permits testing the original lattice
at height-one primes of the normal chart:
a projective lattice is the intersection of its
height-one localizations inside its generic vector
space. Faithful valuation extension detects the
same inclusion. At primes not above \(p\),
(IC6) already gives \(\omega'_p=\omega\).
This proves (IC5) on the chart and hence on the model.

Dualizing (IC5), and using the locally free labelled
restriction modules, gives the exact specialization
statement for the Hodge--Tate dual lattice \(D_{p,s}\):
\[
 \Lambda_{\mathfrak p}\cap V_s
       \subseteq (D_{p,s})_z,\qquad
 \kappa_p(D_{p,s})_z
       \subseteq\Lambda_{\mathfrak p}\cap V_s
                                      \quad(\mathfrak p\mid p).
 \tag{IC7}
\]
To justify base change here, the map from the locally
free \(\omega'_p\) onto its labelled image has
projective quotient and therefore splits locally.
Its dual and the labelled image commute with the
point's base change. Thus (IC7) does not interchange
tensor product with a nonsaturated intersection.
For completeness this specialization also works when
the point is non-dominant. Locally let \(P=\omega\),
\(H=\omega'_p\), \(n=p^{R_p}\), and let
\(I_s=\operatorname{im}(H\to V_s^\vee)\).
The model hypotheses make \(H,I_s\) locally free,
and \(H\to I_s\) is split. There are maps
\(H\to P\) and \(P\to H\) whose composites are
multiplication by \(n\), since \(nP\subset H\subset P\).
For a map to a characteristic-zero valuation domain
\(R\), these identities survive tensoring.
As \(H_R\) is free and \(n\ne0\), they imply that
\(H_R\to P_R\) is still injective.
The split dual \(I_{s,R}^\vee\) embeds in
\(H_R^\vee\subset n^{-1}P_R^\vee\).
Over \(\operatorname{Frac}R\), the labelled
characteristic-zero action identifies this subspace
with the specialized \(s\)-eigenspace.
An integral eigentangent in \(P_R^\vee\) annihilates
the kernel of \(H_R\to I_{s,R}\): this is true over
the fraction field and hence over the domain.
It therefore belongs to \(I_{s,R}^\vee\).
Conversely a member of \(I_{s,R}^\vee\), evaluated
on \(\kappa_pP_R\subset H_R\), gives an integral
functional after multiplication by \(\kappa_p\).
These are exactly the two inclusions in (IC7).
No map of the model's function field to
\(\operatorname{Frac}R\) has been used.
The action and label identification in
\(\mathrm{LinPt}\) make this \(V_s\) the point's
actual eigenspace.

At \(p\), the patched lattice is \(D_{p,s}\) for one
sign of the CM label and its different-saturation
for the other sign. In either case its definition
gives
\[
 D_{p,s}\subseteq\mathfrak M.V_s,\qquad
 \sigma_s(\mathfrak D_{E/\mathbb Q})\,\mathfrak M.V_s
                                      \subseteq D_{p,s}.
 \tag{IC8}
\]
The saturation in the second formula is imposed before
evaluation; multiplication by each \(d\) therefore
preserves that inclusion under every point map.
Combining (IC7)--(IC8) proves the local comparison
\[
 \Lambda_s\otimes\mathcal O_{L,\mathfrak p}
     \subseteq F_{s,\mathfrak p},\qquad
 \kappa_p\,\sigma_s(\mathfrak D_{E/\mathbb Q})
       F_{s,\mathfrak p}
     \subseteq\Lambda_s\otimes\mathcal O_{L,\mathfrak p}.
 \tag{IC9}
\]
Off \(S\), both the patching prime and the different
are units. Equations (IC2)--(IC3), together with the
ordinary Lie/cotangent duality of a smooth group,
give equality of these two local lattices.

\begin{srproposition}\label{prop:sr-actual-integral-eigenlattice}
For the actual compactification, point and
\(\mathrm{HTFree}\) data above, and every CM label \(s\),
there is a nonzero integral ideal \(\mathfrak a_s\)
of \(\mathcal O_L\) such that
\[
             \mathfrak a_s F_s\subseteq\Lambda_s,\qquad
             \mathfrak a_s\Lambda_s\subseteq F_s,
 \tag{IC10}
\]
and
\[
 \frac{\log N_{L/\mathbb Q}\mathfrak a_s}{[L:\mathbb Q]}
       \le
       2\sum_{p\in S}\frac{\log p}{p-1}
                              +2\log\operatorname{rd}(E).
 \tag{IC11}
\]
In fact \(\Lambda_s\subseteq F_s\), and the explicit
ideal below gives the stronger bound with
\(\log\operatorname{rd}(E)\) in place of its twice.
\end{srproposition}
\begin{proof}
The functions \(\kappa_p\) have well-defined values
\(k_p\in L\) at the point. To obtain a value, choose
one prime of \(L\) above \(p\); regularity in (IC5)
gives the value from that local chart.
Uniqueness of evaluation at the generic point makes
this the same value in every local comparison.
Its polynomial relation is preserved by evaluation:
\[
       k_2=4,\qquad
                  \Phi_p(k_p+1)=0\quad(p\ne2).
 \tag{IC12}
\]
Since \(L\) has characteristic zero, \(k_p+1\) is a
primitive \(p\)-th root, and \(k_p\ne0\).
These relations also prove
\(k_p\in\mathcal O_L\), globally.
Set
\[
 K_S=\prod_{p\in S}k_p,\qquad
 \mathfrak a_s=
    \sigma_s(\mathfrak D_{E/\mathbb Q})\mathcal O_L
                                      \cdot(K_S).
 \tag{IC13}
\]
Both factors are nonzero.

The local saturated evaluation defining \(F_s\)
has the following membership criterion:
\(x\in F_s\) if and only if, at every prime
\(\mathfrak p\), some \(t\notin\mathfrak p\)
makes \(tx\) a finite
\(\mathcal O_L\)-linear combination of evaluated
local generators of the patched lattice.
This is just localization of a span.
Thus (IC9) gives
\(\Lambda_s\subseteq F_s\) globally.
For \(x\in F_s\) and
\(d\in\mathfrak D_{E/\mathbb Q}\), it also gives
\[
                           K_S\sigma_s(d)x\in V_s
\]
and says that this vector is an integral tangent
vector over every \(\mathcal O_{L,\mathfrak p}\).
The other \(k_q\)'s are integral and preserve
these inclusions.

The integral tangent module is finite projective,
since the group is smooth.
Over the Dedekind ring \(\mathcal O_L\) it is the
intersection of its localizations inside \(V\).
Hence the vector just displayed belongs to
\(\Lambda_s\).
Equivalently, write its coordinates in local
trivializations of the tangent vector bundle:
membership in every DVR localization makes every
coordinate integral, and the unique local tangent
sections agree on the fraction field and glue.
This is exactly the
\(\mathrm{tangent\_vector\_extends\_dvr}\)
step, with the same cartesian Lie identification
as in (IC1).
Taking the ideal generated by all these \(d\)
proves \(\mathfrak a_sF_s\subseteq\Lambda_s\).
The other inclusion in (IC10) follows from
\(\Lambda_s\subseteq F_s\) and integrality of
\(\mathfrak a_s\).

For odd \(p\), the cyclotomic identity gives
\[
                 (k_p)^{p-1}=(p)
                    \quad\text{as ideals of }\mathcal O_L .
\]
Thus
\[
 \frac{\log N(k_p)}{[L:\mathbb Q]}
               =\frac{\log p}{p-1}\quad(p\ne2),\qquad
 \frac{\log N(4)}{[L:\mathbb Q]}=2\log2 .
 \tag{IC14}
\]
The latter is exactly twice
\(\log2/(2-1)\); no replacement \(k_2=2\) is made.

For the extended different, the tower norm formula is
\[
 N_{L/\mathbb Q}
  \bigl(\sigma_s(\mathfrak D_{E/\mathbb Q})\mathcal O_L\bigr)
          =|\operatorname{disc}E|^{[L:E]},
\]
where \(L\) is regarded as an \(E\)-field through
\(\sigma_s\). Therefore its normalized logarithm is
\(\log\operatorname{rd}(E)\).
Combining this with (IC13)--(IC14) yields
\[
 \frac{\log N\mathfrak a_s}{[L:\mathbb Q]}
       =\log\operatorname{rd}(E)
          +2\,\mathbf1_{2\in S}\log2
          +\sum_{\substack{p\in S\\p\ne2}}
                                    \frac{\log p}{p-1}.
 \tag{IC15}
\]
Since \(\operatorname{rd}(E)\ge1\), this implies
(IC11), including all ramification and field-degree
normalizations. Finite extension of the point field
multiplies both logarithmic ideal norms and its
degree by the same factor.
\end{proof}

\paragraph{Source and hypothesis alignment.}
The exact final producer is
\texttt{Development/WeilBound/CoreStatements/%
PlurisubharmonicCones/AbelianPartIntegralPointBis.lean},
theorem
\texttt{LevelTowerAlt.abelian\_part\_integral\_point\_bis}.
Its exported statement is
\texttt{GeometryProofs/AbelianPartIntegralPoint.lean}.
The local calculation is
\texttt{LocalPrimeBound.on\_S}, \texttt{off\_S},
and \texttt{loc}, in
\texttt{AevalCyclotomic.lean} and \texttt{InFibreLoc.lean}.
The explicit ideal and norm are \texttt{NormBound.lean};
the different norm is
\texttt{PDifferentExtNorm.absNorm\_extDifferent}.
The finite-level family and test-object joins are
\texttt{HodgeTateFreeFlagLatticeLower},
\texttt{HodgeTateFreeFlagLatticeUpper},
\texttt{PointHTOfFarguesHodgeTate}, and
\texttt{ExistsSandwichOfHodgeTateFree}.
The two freeness clauses above are used at the
projective-duality and specialization steps.
They are not consequences of the Fargues bound alone.

\subsection{The initial torsion level and the freeness modification}
\label{subsec:sr-initial-htfree-construction}

We record how the freeness hypothesis is produced in the
actual root-model construction, separating it from later
transport of already free models.

Let \(e_0\) be the chosen generic component of the original
integral moduli model and fix \(S\) and \(R_p\ge4\).
Set \(n=\prod_{p\in S}p^{R_p}\).
The generic fibre has characteristic zero, so its
\(n\)-torsion and its full symplectic level structures
form a finite etale cover.
Choose one component of this cover and normalize the
projective closure of \(e_0\) in its finite function-field
extension. Normalization is finite here because the
model is excellent. This gives a finite normal
projective model with full generic level \(n\).
The principal polarization and a symplectic pair with
pairing \(1\) give a primitive \(n\)-th root of unity
in this function field, hence all required
\(p^{R_p}\)-th roots.
The chosen component is compatible with the integral
PEL Tate parametrization, including the
\(\mathcal O_E\)-action and the alternating form.
For precision, choose over an algebraic closure of
the generic field the Tate parametrization supplied
by the generic moduli datum, descend its finitely many
values on \(\mathcal L/n\mathcal L\) and its primitive
root, and take the resulting finite level field.
Every embedding of this field into an algebraically
closed field extends to the original algebraic
closure. Transport of the chosen parametrization
then gives a compatible full Tate parametrization
extending the finite level over every such field.
This is exactly
\texttt{IsFullLevelAt}, as produced by
\texttt{ExistsFullLevelFiniteExtension.exists\_fullLevel\_finite};
it is stronger data than merely choosing a
symplectic basis of one finite torsion group.
Pull back the given semiabelian group, action and
cubical polarization.

On a normal affine open, the quasi-finite
\(p^{R_p}\)-torsion model and its Weil characters
give the representatives in (IC6).
The character values, being roots of unity, are
integral units on the normal base.
A local lift of their differential gives \(h_x\);
two lifts differ by \(p^{R_p}\omega\).
There are finitely many generic torsion points,
so (IC6) is a finite generating set.
It proves coherence of \(\omega'_p\), full generic
rank, and equality with \(\omega\) after inverting \(p\).
The images in every labelled generic dual space
are coherent as images of finite modules, with
the prescribed generic ranks.
No integral decomposition of the \(\mathcal O_E\)
action at ramified primes is used to form these
images.

Here is the local existence argument for (IC6),
including a boundary point.
Pass first to the normal henselian local base at
that point. Multiplication by \(n=p^{R_p}\) on the
semiabelian group has quasi-finite flat kernel.
Its finite open part containing the entire closed
fibre is a finite flat subgroup \(H\).
This finite part is obtained by Zariski's main
theorem and the finite/open decomposition over a
henselian base; its uniqueness makes it a subgroup.

Full generic level identifies \(H_\Omega\) with a
subgroup of the fixed finite torsion group.
For a generic torsion point \(x\), its polarized
Weil character restricted to \(H_\Omega\) is a
generic section of the finite Cartier dual \(H^D\).
A section of a finite scheme over a normal domain
extends from its fraction field: every coordinate
of the section is integral over the base and
therefore belongs to it. Thus this character
extends over the henselian local base.
Near the identity, \(H\) is the full kernel of
multiplication by \(n\).
Lift the character function to the smooth ambient
group and differentiate at the identity.
Two lifts differ by the kernel's conormal ideal;
its image in the ambient cotangent module is
\[
                         d[n]^*\omega=n\omega .
\]
This proves the asserted ambiguity.
This construction gives a class modulo \(n\omega\)
on an elementary etale neighbourhood: the finite
part, its character and the chosen function descend
there by finite presentation.
On a normal etale overlap compare the two finite
open parts near the identity.
Their intersection is finite flat, open and closed
in each, and has dense generic fibre.
The two characters restrict to the same generic
Weil character, hence agree on this intersection.
Their differentiated classes therefore agree modulo
\(n\omega\), since near the identity its defining
ideal is the same kernel-of-\([n]\) ideal.
Etale descent for the finite module
\(\omega/n\omega\) now gives a unique class on the
original local base.
A local lift in the finite free cotangent module
gives the required \(h_x\).
For the finite-part decomposition one can use
\href{https://stacks.math.columbia.edu/tag/02LN}
{Stacks Project, Lemma 37.41.4}, applying it to all
points of the finite closed fibre, then henselizing;
the flatness comes from that of the torsion kernel.
This constructs the local forms without asserting
that \(x\) itself extends to a boundary section.

The pullback assertion here concerns dominant morphisms
between normal models carrying the transported generic
torsion labels and cubical pairing. These include the
Grassmann modifications and purely transcendental flag
extensions used below. The pulled-back finite part is
finite open and contains the entire new closed fibre;
finite-clopen uniqueness identifies it with the new
finite part. Its characters have the same generic
values and therefore agree by normality.
This proves equality of the generated differential
lattices, including on these nonflat modifications.
It does not identify the intrinsic Hodge--Tate lattice
of an arbitrary specialized semiabelian fibre with
a pullback of the generic abelian torsion construction.
At an arbitrary integral point the target uses the
evaluated coherent model lattice, as in (IC1).
These are the constructions expanded in
\texttt{HodgeTateWeilLocal.exists\_tors},
\texttt{exists\_form}, \texttt{htLat\_exists\_lifts},
and \texttt{HodgeTateWeilPullback.main}.

There is one projective modification making this
finite list of coherent full-rank lattices locally free.
For a coherent torsion-free module of generic rank \(r\),
take its relative Grassmannian of rank-\(r\) quotients,
close the generic identity quotient, and normalize.
The universal quotient is locally free.
Its kernel is generically zero and hence torsion;
inside the fixed generic vector space the quotient
is exactly the image of the pulled-back lattice.
For several modules take the closure of the common
generic section in the product of these Grassmannians
before normalizing.
Apply this to every \(\omega'_p\) and every one of
its labelled dual restrictions, for every \(p\in S\).
The finite family is essential.

The pulled-back lattices are the actual finite-level
lattices of the new model.
Indeed pull back the representatives \(h_x\) and the
ordinary cotangent basis; they retain the same Weil
characters. Every new representative for a torsion
point differs from its corresponding pulled-back one
by \(p^{R_p}\omega\).
Both (IC6) and this ambiguity therefore give equality
of the generated modules in both directions.
The same equality for their labelled images is
right-exactness of image and evaluation inside
the fixed generic labelled space.
Thus these are the actual \(\mathrm{HTFree}\)
lattices, not unrelated flattened replacements.

On the rational generic fibre all these modules
were already vector bundles: \(p\) and the different
are units there, and the labelled action is etale.
The Grassmann modification is therefore an isomorphism
there. It changes neither the generic component nor
its finite cover degree.
The different-twisted eigenspace lattices and their
duals are locally free as well, since the different
is an invertible ideal of the Dedekind ring
\(\mathcal O_E\).
For completeness the subsequent flagged root model
can be built without a new freeness assertion.
Put \(r_s=\operatorname{labelRk}(D,s)\).
Take the generic point of the product of the complete
flag varieties of these labelled generic spaces.
In its big cell its field is the original field
with \(\sum_s r_s(r_s-1)/2\) independent parameters
adjoined. Use a basis over the original unlevelled
generic field to choose these parameters, so a finite
level-cover degree \(d\) remains \(d\) after this
purely transcendental extension:
\([K(\mathbf t):K_0(\mathbf t)]=[K:K_0]=d\).
This is the \(\mathrm{IsGenericFlag}\) degree in
\texttt{HodgeCptGoodReductionDegree}; the file's name
does not introduce an additional good-reduction
hypothesis at this step.

Take the normal projective closure of the flag graph.
A common flag/Grassmann modification, applied to the
finite list of flag-step quotients both in the
patched labelled lattices and in the full Lie
lattice, makes their strict transforms locally free.
Consequently every flag step and every successive
graded piece is a vector bundle of the prescribed
rank. Pullback of the actual Hodge--Tate lattices is
the equality already proved, and local freeness
of the previously constructed \(\omega'_p\) and
its labelled images is preserved.
Thus both \(\mathrm{HTFree}\) clauses survive these
modifications; freeness of a flag alone is not
substituted for them.

The root lattices in the source are the specified
tensor powers, including negative powers, of these
rank-one graded pieces.
They are therefore actual line bundles.
For each one, choose a sufficiently high power of
an ample coordinate bundle making its tensor
product globally generated.
Adding the resulting finite coordinate block
expresses that line bundle as an integral
combination of projective-coordinate bundles,
with its chosen rational generator.
Do this for the finite set of all root lattices,
retain the old coordinate blocks, and use a common
Segre refinement.
This produces exactly the integers \(x_{\rm deg}\)
and functions \(e\) in
\(\mathrm{Fl.RootLines}(x_{\rm deg},e)\), on the
same normal projective model.
No extraction of new algebraic roots, or new
field-degree factor, is part of this coordinate
step.
This is the construction in
\texttt{HodgeCptLabelBases.exists\_flag\_model}
and \texttt{HodgeCptBlockRefinement}.
The parameter count also gives
\[
 \dim(\text{PEL generic fibre})
   +\sum_s\frac{r_s(r_s-1)}2
       =f\,\frac{(p+q)(p+q-1)}2,
\]
because the two ranks over each CM place are \(p,q\)
and the PEL dimension there is \(pq\).
It is the arithmetic dimension required by the
root model after adding the base dimension.

This is the initial construction in
\texttt{HodgeCptProjectiveModel.exists\_hodgeCpt\_htFree},
preceded by
\texttt{LevelCoverStHolds.levelCover}
and assembled in
\texttt{IsPrimitiveRootOfLevel.exists\_hodgeCpt\_flag}.
The later
\texttt{HodgeCompactCoverPrimes.hodge\_compact\_cover\_primes}
preserves this already established hypothesis while
adjoining the root model; it is not used as its
initial producer.
The construction assumes the actual polarized
semiabelian model and its finite-level Weil
uniformization data. Those are the common geometric
background of the moduli model, independent of the
Fargues theorem and of this lattice comparison.
